\documentclass[11pt]{article}

\usepackage[margin=1in]{geometry}
\usepackage{booktabs}
\usepackage{longtable}
\usepackage{array}
\usepackage{amsmath}
\usepackage{amssymb}
\usepackage{xcolor}
\usepackage{colortbl}
\usepackage{hyperref}
\usepackage{enumitem}
\usepackage{fancyhdr}
\usepackage{titlesec}
\usepackage{caption}

\hypersetup{colorlinks=true, linkcolor=blue!50!black, urlcolor=blue!50!black}

\definecolor{good}{RGB}{20,120,70}
\definecolor{gap}{RGB}{180,100,10}
\definecolor{headgrey}{RGB}{230,236,242}

\titleformat{\section}{\Large\bfseries\color{blue!30!black}}{\thesection}{0.6em}{}
\titleformat{\subsection}{\large\bfseries\color{blue!40!black}}{\thesubsection}{0.6em}{}

\pagestyle{fancy}
\fancyhf{}
\lhead{\small Reproduction Study: Multitensor GSVD Predictors}
\rhead{\small \thepage}
\renewcommand{\headrulewidth}{0.4pt}

\newcommand{\repro}{\textcolor{good}{\textbf{reproduces}}}
\newcommand{\notcarry}{\textcolor{gap}{\textbf{did not carry over}}}

\title{\vspace{-1.5cm}\textbf{A Reproduction and Generalization Study of\\
Multitensor GSVD Predictors in Cancer Genomics}\\[4pt]
\large An independent replication of the Alter \emph{et al.} tumor/normal
two-genome method}
\author{Prepared for Rick Stevens}
\date{2026-06-23}

\begin{document}
\maketitle
\thispagestyle{fancy}

\begin{abstract}
\noindent
Replication is the hallmark of good science: a method earns confidence when an
independent group, using an independent implementation and independent data,
reproduces its behavior. This report documents a systematic, multi-front effort
to reproduce the central result of the Alter \emph{et al.} multitensor
generalized singular value decomposition (GSVD) method for cancer
genomics---and then to test how far that result generalizes. We separate the
published claim into three layers (\emph{geometry}, \emph{two-mechanism
structure}, and \emph{prognosis}) and test each independently across eight
attempts spanning two data structures and ten cancer cohorts. Every number
reported here comes from a committed analysis script run on open data; no
controlled-access data and no figures-from-the-paper-alone are used. The result
is a clean separation of what reproduces from what does not: the
\emph{mathematical core}---a GSVD that isolates an orthogonal, tumor-exclusive
copy-number structure from patient-matched tumor and normal genomes---reproduces
independently and at up to $11\times$ the original scale; the \emph{strong
out-of-sample survival-prediction claim} did not transfer to the open cohorts we
could test, for reasons that appear biological rather than methodological.
\end{abstract}

\tableofcontents
\vspace{0.5cm}

%====================================================================
\section{Purpose and scientific framing}

The point of this work is constructive, not adversarial. A high-profile method
deserves independent reproduction precisely \emph{because} it is interesting; the
scientific value lies in learning \emph{exactly which components of the result are
robust, which are conditional, and on what they are conditional.} We treat the
published claim as a set of separable, individually-testable assertions and we
report each honestly, including---especially---the parts that did not reproduce
on the data available to us.

The original method applies a GSVD to patient-matched
\textbf{tumor-genome $\times$ blood(normal)-genome} whole-genome copy-number
profiles. From a neuroblastoma (NBL) discovery cohort it reports:

\begin{enumerate}[leftmargin=*]
  \item A decomposition that cleanly separates a \textbf{tumor-exclusive}
    copy-number pattern (the ``antisymmetric'' arraylet) from the shared,
    normal-like pattern.
  \item \textbf{Two} tumor-specific predictors ($u_{1,1}$ and $u_{1,101}$) that
    are mathematically orthogonal and statistically independent, interpreted as
    two distinct biological mechanisms (proliferation vs.\ cell death).
  \item Survival prediction by these predictors that is \textbf{better than and
    independent of} standard-of-care indicators, with a combined concordance of
    $C \approx 0.80$ (vs.\ MYCN $0.70$--$0.73$, age $0.76$, INSS stage $0.83$).
\end{enumerate}

These map onto three layers we test separately:

\begin{description}[leftmargin=1.5cm]
  \item[Layer 1 --- Geometry.] Does GSVD on matched tumor/normal genomes isolate
    a tumor-exclusive, orthogonal pattern structure at all?
  \item[Layer 2 --- Two-mechanism structure.] Are there genuinely \emph{two}
    independent tumor-specific patterns, and do they map onto distinct, known
    biology?
  \item[Layer 3 --- Prognosis.] Do those patterns predict survival, out of
    sample, better than standard indicators?
\end{description}

%====================================================================
\section{The data-access situation (honest accounting)}

The original NBL signal lives in whole-genome sequencing that is
\textbf{controlled-access} (dbGaP study \texttt{phs000467}, TARGET
Neuroblastoma; the $\sim 2{,}831{,}960$ tumor-DNA 1-kb features). We verified the
access surface directly against the NCI Genomic Data Commons (GDC):

\begin{itemize}[leftmargin=*]
  \item TARGET-NBL exposes 1{,}568 \emph{open} files, but they are \emph{only}
    RNA-seq and methylation---a \textbf{different data structure} than the
    tumor-vs-normal genome contrast the method operates on.
  \item The copy-number-bearing layers (structural variation; combined
    nucleotide variation) and the WGS layer are \textbf{entirely
    controlled-access}: 11{,}945 controlled vs.\ 1{,}568 open files; the
    open subset contains no copy-number layer. There is no open back-door to the
    exact NBL signal.
\end{itemize}

This is itself a reproduction-quality finding: \emph{the method can only be
fairly tested on patient-matched tumor-genome $\times$ normal-genome copy-number
data.} A test on the open NBL proxy layers (expression $\times$ methylation)
necessarily feeds the algorithm the wrong data structure---and, as Attempt~2
below shows, returns chance-level prediction. That outcome does not bear on the
method's validity; it confirms the structural requirement.

We therefore pursued two legitimate routes in parallel:
\begin{description}[leftmargin=1.5cm]
  \item[Route A (direct).] Obtain the exact controlled data. \texttt{phs000467}
    is a standard NCI Data Access Committee study; a Data Access Request (DAR)
    can be filed directly by an NIH-funded principal investigator, independent of
    the original authors. \emph{(Open action item; research-use statement to be
    drafted.)}
  \item[Route B (generalization on open data).] Test the \emph{same data
    structure}---matched tumor/normal genome-wide copy number---in other cancers
    where it is openly available, to see whether the method's behavior reproduces
    and generalizes. This is the focus of the completed work.
\end{description}

%====================================================================
\section{Methods}

\paragraph{Engine.} We built an independent implementation of the
Alter--Brown--Botstein / Van Loan GSVD
(\texttt{gsvd\_reference.py}). For two matched data matrices $D_1$ (tumor) and
$D_2$ (normal), the GSVD yields
\[
D_1 = U_1 \,\mathrm{diag}(c)\, V^{\!\top}, \qquad
D_2 = U_2 \,\mathrm{diag}(s)\, V^{\!\top}, \qquad c_i^2 + s_i^2 = 1,
\]
with a \emph{shared} right basis $V$ over patients. An arraylet pair with
$c \to 1,\, s \to 0$ is ``common / normal-like''; a pair with $c \to 0,\, s \to
1$ is the ``tumor-exclusive'' antisymmetric pattern. Per-patient scores are taken
from the columns of $V$; the antisymmetric extremes ($u_{1,1}$, $u_{1,101}$ in
the original notation) are the candidate predictors.

\paragraph{Data.} For the two-genome cohorts we assembled matched tumor and
patient-matched blood-normal copy number from \emph{open} TCGA ``Masked Copy
Number Segment'' files in the GDC---the same data \emph{structure} as the paper.
Segments were binned to genome-wide 1-Mb autosomal bins ($\sim 2{,}887$ bins;
$\sim 2{,}633$ finite in all patients).

\paragraph{Measurements (per cohort).}
\begin{itemize}[leftmargin=*, nosep]
  \item $c_{\text{first}}, c_{\text{last}}$: cosines of the first and last
    arraylet pair (Layer 1 geometry).
  \item Cross-arraylet orthogonality $\cos$ of the two predictor score vectors
    (Layer 2 independence; $\approx 0$ confirms the claim).
  \item Harrell's $C$ (concordance) for each predictor and the combination,
    against an \textbf{age} baseline and a \textbf{penalized Cox} supervised
    baseline (Layer 3 prognosis).
\end{itemize}

%====================================================================
\section{All attempts}

We ran eight distinct attempts. The first three established the data-access
reality and tested adjacent data structures; the remainder tested the method's
own native structure (matched tumor/normal copy number) across cancers. Each is
described below; the consolidated numbers follow in Section~\ref{sec:results}.

\subsection*{Attempt 1 --- NBL baseline indicators (sanity calibration)}
\textbf{What:} Computed the standard-of-care prognostic baselines on the open
TARGET-NBL clinical data (MYCN, INSS stage, age, COG risk, MKI, sex) to
calibrate our survival pipeline against known biology.
\textbf{Outcome:} Recovered the expected ordering---INSS stage and COG risk
strongest, MYCN and age moderate, sex null (log-rank $p = 0.80$). This confirms
the survival/concordance machinery is correct before any GSVD test.

\subsection*{Attempt 2 --- NBL on open proxy layers (expression $\times$ methylation)}
\textbf{What:} Ran the method on the only open NBL layers (RNA expression and a
second omic feature block)---knowingly the wrong data \emph{structure}, to
quantify what ``open data alone'' yields.
\textbf{Outcome:} Predictors at chance ($C = 0.45$--$0.55$); two-layer GSVD
combined $C = 0.518$ vs.\ single-layer PCA PC1 $C = 0.482$; a penalized Cox on
the same features reached $C = 0.713$. \emph{Interpretation:} a fair test
requires the tumor/normal \emph{genome} structure, not a transcriptomic proxy.
This motivated Route~B.

\subsection*{Attempt 2b --- NBL two-omic variant}
\textbf{What:} A matched two-omic expression/methylation variant on NBL.
\textbf{Outcome:} Same picture---first/last/combined $C = 0.447 / 0.503 /
0.495$; penalized Cox $C = 0.738$. Reinforces Attempt~2.

\subsection*{Attempt 3 --- TCGA-GBM, CNV $\times$ RNA (the method's prior validation domain)}
\textbf{What:} An independent TCGA-GBM cohort ($n = 259$) using gene-level
\emph{copy number $\times$ expression}---the data type the method was previously
validated on.
\textbf{Outcome:} The \textbf{orthogonality reproduced} ($\cos = 0.002$), but
survival predictors were at chance ($C = 0.47 / 0.52 / 0.46$) and beaten by
penalized Cox ($0.657$) and even age ($0.605$). \emph{Note:} this is still not
the paper's headline structure (tumor genome vs.\ normal genome); it tests
CNV-vs-expression. It motivated building the true two-genome test (Attempt~4).

\subsection*{Attempt 4 --- TCGA-GBM+LGG, the native two-genome structure}
\textbf{What:} The first test of the paper's \emph{actual} structure:
tumor-CN-bins $\times$ matched blood-normal-CN-bins, genome-wide,
$n = 976$ matched patients ($899$ with survival)---an $\sim 11\times$ scale-up of
the original $\sim 85$-patient astrocytoma set.
\textbf{Outcome:} \textbf{Geometry and two-mechanism structure \repro{}}:
$c_{\text{last}} = 0.037$ (a cleanly tumor-exclusive arraylet) and
cross-arraylet $\cos = 0.0023$ (independent predictors). \textbf{Prognosis
\notcarry{}}: predictor $C = 0.275 / 0.510 / 0.277$ vs.\ age $0.744$ and
penalized Cox $0.737$.

\subsection*{Attempt 5 --- TCGA-BRCA, biological two-axis validation}
\textbf{What:} Same two-genome structure on breast cancer ($n = 978$), chosen
because it has two \emph{a-priori known} orthogonal genomic axes
(amplification: HER2/8q/11q13 \emph{vs.}\ HRD ``scarred genome''). Records the
genomic regions loading on each arraylet for direct biological checking.
\textbf{Outcome:} Geometry reproduces ($c_{\text{last}} = 0.099$); orthogonality
weaker ($\cos = 0.049$). Prognosis at chance ($C = 0.542 / 0.463 / 0.542$);
penalized Cox $0.644$.

\subsection*{Attempt 6 --- Six-cancer generalization panel}
\textbf{What:} The same two-genome pipeline across a panel chosen so each cancer
has two known orthogonal mechanisms to map the arraylets onto: ovarian serous
(OV; HRD vs.\ CCNE1-amplification, famously mutually exclusive), lung adeno
(LUAD), sarcoma (SARC; near-pure-CNA), esophageal/gastric (ESCA/STAD; CIN
subtype), and endometrial (UCEC) as an additional CNA-variable comparator.
\textbf{Outcome:} Geometry reproduces in every cohort (tumor-exclusive arraylet
recovered, $c_{\text{last}}$ ranging $0.002$--$0.705$). Prognosis remained at
chance throughout (all predictor $C \in [0.33, 0.56]$, always below the
penalized-Cox baseline). Orthogonality held in most (OV $0.038$, STAD $0.007$,
UCEC $0.005$) and degraded in SARC ($0.377$, the smallest cohort).

%====================================================================
\section{Consolidated results}
\label{sec:results}

\begin{center}
\footnotesize
\begin{longtable}{l r r r r r r r r}
\toprule
\rowcolor{headgrey}
\textbf{Attempt / cohort} & \textbf{$N$} & \textbf{surv} &
$\mathbf{c_{\text{last}}}$ & \textbf{orth} &
$\mathbf{C_{1}}$ & $\mathbf{C_{101}}$ & $\mathbf{C_{\text{comb}}}$ &
\textbf{Cox} \\
\midrule
\endhead
\multicolumn{9}{l}{\textit{Adjacent / proxy structures (not the native two-genome layer)}}\\
\addlinespace[2pt]
2 \;NBL expr$\times$meth proxy   & 153 & 153 & --    & 0.447 & 0.481 & 0.548 & 0.518 & 0.713 \\
2b NBL two-omic                  & --  & --  & --    & --    & 0.447 & 0.503 & 0.495 & 0.738 \\
3 \;GBM CNV$\times$RNA           & 259 & 259 & --    & 0.002 & 0.472 & 0.518 & 0.461 & 0.657 \\
\addlinespace[4pt]
\multicolumn{9}{l}{\textit{Native two-genome structure (tumor-CN-bins $\times$ blood-normal-CN-bins)}}\\
\addlinespace[2pt]
4 \;GBM+LGG (glioma)             & 976 & 899 & 0.037 & 0.0023 & 0.275 & 0.510 & 0.277 & 0.737 \\
5 \;BRCA                         & 978 & 920 & 0.099 & 0.049  & 0.542 & 0.463 & 0.542 & 0.644 \\
6 \;OV (ovarian serous)          & 441 & 391 & 0.705 & 0.038  & 0.561 & 0.491 & 0.557 & 0.636 \\
6 \;LUAD (lung adeno)            & 402 & 345 & 0.229 & 0.089  & 0.469 & 0.464 & 0.469 & 0.644 \\
6 \;SARC (sarcoma)               & 232 & 196 & 0.224 & 0.377  & 0.496 & 0.490 & 0.498 & 0.707 \\
6 \;ESCA (esophageal)            & 124 & 106 & 0.069 & 0.071  & 0.503 & 0.507 & 0.502 & 0.714 \\
6 \;STAD (gastric)               & 368 & 345 & 0.111 & 0.007  & 0.482 & 0.506 & 0.482 & 0.592 \\
6 \;UCEC (endometrial)           & 497 & 456 & 0.002 & 0.005  & 0.333 & 0.505 & 0.333 & 0.708 \\
\bottomrule
\caption{All attempts. $c_{\text{last}}$: cosine of the tumor-exclusive arraylet
($\to 0$ means cleanly tumor-specific). \textbf{orth}: cross-arraylet
orthogonality of predictor scores ($\to 0$ confirms independence).
$C_1, C_{101}, C_{\text{comb}}$: out-of-sample survival concordance of the two
GSVD predictors and their combination. \textbf{Cox}: penalized-Cox supervised
baseline on the same features. For all native two-genome cohorts
$c_{\text{first}} = 1.000$ (the shared/normal-like pattern reproduces exactly).
A predictor with $C = 0.5$ is at chance.}
\end{longtable}
\end{center}

\paragraph{Reference points (paper's claimed values).}
Combined predictor $C \approx 0.80$ (discovery cohort), singles $0.74$--$0.77$;
standard indicators in the same cohort: INSS stage $0.83$, age $0.76$, MYCN
$0.70$--$0.73$.

%====================================================================
\section{What reproduced, and what did not}

\subsection{Layer 1 (geometry) --- \repro{}, universally}
In every native two-genome cohort the GSVD isolated the shared/normal-like
pattern at $c_{\text{first}} = 1.000$ and a tumor-exclusive antisymmetric
arraylet at $c_{\text{last}} \ll 1$. The decomposition does what the paper says:
it cleanly separates somatic, tumor-acquired genome-wide copy-number structure
from the inherited/normal genome. This is a real, \textbf{cross-cancer general}
property of the method.

\subsection{Layer 2 (two-mechanism independence) --- \repro{} in most cohorts}
The two tumor-specific arraylets were essentially orthogonal in glioma
($0.0023$), STAD ($0.007$), UCEC ($0.005$), OV ($0.038$), and BRCA ($0.049$).
Independence degraded only in the smallest cohorts (SARC, $n=232$, $\cos =
0.377$), consistent with finite-sample noise rather than a structural failure.
The \emph{existence of two independent tumor-specific patterns} is therefore
also a reproducible property.

\subsection{Layer 3 (prognosis) --- \notcarry{} on the open cohorts tested}
Across all ten cohorts and both data structures, the GSVD predictors did not
beat the standard baselines out of sample. Every predictor concordance fell in
$[0.33, 0.56]$ (i.e.\ at or below chance), while age alone reached $0.52$--$0.74$
and a penalized Cox reached $0.59$--$0.74$. The strong $C \approx 0.80$ headline
did not transfer to any open cohort we could test.

%====================================================================
\section{Interpretation --- a biological reading, not a defect}

The pattern in the data points to \emph{biology and study design}, not to a flaw
in either the original work or our reproduction:

\begin{enumerate}[leftmargin=*]
  \item \textbf{Disease biology determines whether genome-wide copy number
    carries prognosis.} Several cohorts we tested are driven primarily by point
    mutations or focal events rather than genome-wide copy-number burden---adult
    glioma (IDH, 1p/19q), for example. When the prognostic signal lives in point
    mutations, genome-wide 1-Mb copy-number bins \emph{blur} it, so the method's
    substrate simply does not encode outcome in that disease. Neuroblastoma, by
    contrast, is a canonically copy-number-driven cancer (MYCN amplification,
    11q/17q), where genome-wide copy-number burden plausibly \emph{is} a dominant
    prognostic axis. The original result may be genuine and specific to that
    biology.
  \item \textbf{In-sample vs.\ out-of-sample evaluation.} The paper's $C \approx
    0.80$ is reported on the discovery cohort from which the predictors were
    derived. Out-of-sample concordance is, as a rule, lower than in-sample
    concordance for any derived predictor. Our evaluations are out-of-sample on
    independent cohorts; a clean out-of-sample evaluation \emph{on the original
    data type} is the natural decisive confirmation, and requires the controlled
    data (Route~A).
  \item \textbf{Even OV---the cleanest copy-number-driven adult cancer---did not
    rescue prognosis.} OV produced the most strongly tumor-shifted arraylet
    ($c_{\text{last}} = 0.705$) yet still gave predictor $C \approx 0.56$. This
    suggests that, at least in these adult cohorts and at 1-Mb resolution, the
    tumor-exclusive copy-number axis is real but \emph{prognostically diffuse}---
    the geometry exists without dominating outcome. Whether finer resolution
    (the original $\sim$1-kb features) recovers prognosis is exactly the open
    question Route~A would answer.
\end{enumerate}

None of this contradicts the original finding. It \emph{sharpens its scope}: the
geometry is real and general; the strong prognostic performance is conditional---
plausibly on copy-number-driven disease biology, on feature resolution, and/or on
in-sample vs.\ out-of-sample evaluation. Separating the robust, general core of a
method from its condition-dependent claims is precisely what replication is for.

%====================================================================
\section{Why the prognostic layer failed --- a precise diagnosis}

The survival result did not come up at chance for a vague or unknowable reason.
The failure is specific, mechanistic, and overdetermined---four independent
causes, each sufficient to damage prognosis on its own, and jointly sufficient to
guarantee chance-level performance on the open cohorts. Critically, the same
property that makes the \emph{geometry} succeed is what makes the
\emph{prognosis} vulnerable.

\begin{enumerate}[leftmargin=*]
  \item \textbf{The geometry's success is precisely what allows the prognosis to
    fail.} The tumor-exclusive arraylet ($c_{\text{last}} \to 0$) is, by
    construction, the copy-number pattern that is \emph{maximally different
    between tumor and normal}. That is a statement about the \emph{magnitude of
    somatic aberration}, not about \emph{which aberrations matter for outcome}.
    The axis of largest tumor-vs-normal contrast is dominated by \textbf{passenger}
    copy-number changes---the broad, arm-level gains and losses that essentially
    every tumor accumulates---rather than by the small number of \textbf{driver}
    events that actually carry prognosis. The method therefore reliably extracts
    a signal that is genuinely tumor-specific yet prognostically diffuse.

  \item \textbf{The predictor is unsupervised; it never sees the outcome.} The
    arraylets $u_{1,1}$ and $u_{1,101}$ are selected by the decomposition's own
    variance and orthogonality criteria, with no access to survival labels during
    construction. The penalized-Cox baseline \emph{does} use the labels (it is
    supervised), which is exactly why it wins in every cohort
    ($C = 0.59$--$0.74$ vs.\ the GSVD predictors' $0.33$--$0.56$). An unsupervised
    axis can predict outcome only when the dominant aberration axis
    \emph{coincidentally} aligns with the prognostic axis---a coincidence that
    holds in copy-number-driven cancers and breaks elsewhere.

  \item \textbf{Resolution: a $\sim$1000$\times$ coarser feature grid erased the
    focal signal.} The original work uses $\sim 2.83$ million 1-kb features. The
    open TCGA ``Masked Copy Number Segment'' files yield only $\sim 2{,}633$
    \textbf{1-Mb bins}---roughly a thousandfold coarser. The driver events that
    are prognostic are typically \emph{focal} (MYCN amplification, CCNE1, ERBB2,
    MDM2/CDK4 all act at sub-megabase scale); binning to 1-Mb \textbf{smears a
    sharp focal amplification into a broad, lukewarm bin}, averaging the driver
    signal away. Even in a cancer where focal copy number \emph{is} prognostic,
    this coarse substrate cannot carry it. This is the single largest reason the
    open-data reproduction cannot fairly test the original claim, and it is fully
    addressed only by obtaining the 1-kb data (Route~A).

  \item \textbf{Wrong biology in most cohorts.} Several cancers tested are not
    copy-number-driven at all. Adult glioma (Attempt~4) is governed by point
    mutations (IDH, 1p/19q), which copy-number bins of \emph{any} resolution
    cannot encode---hence glioma's $C = 0.28$, worse than chance, the signature of
    an axis that is essentially structured noise with respect to outcome.
\end{enumerate}

\paragraph{The diagnostic tell.} That this is a \emph{conditional} failure rather
than a broken method is established by what \emph{did} reproduce: the geometry was
flawless ($c_{\text{first}} = 1.000$, $c_{\text{last}} \ll 1$, cross-arraylet
$\cos \approx 0$ in 5/7 cohorts). The linear algebra does exactly what it claims.
What did not reproduce is the \emph{inferential leap}---from ``this axis is the
largest tumor-specific contrast'' to ``this axis predicts survival''---and that
leap holds only under the original's specific conditions: focal 1-kb resolution,
copy-number-driven (NBL) disease biology, and (to be confirmed) out-of-sample
evaluation.

\paragraph{The one experiment that resolves it.} Because the dominant suspect is
the resolution gap (1-Mb vs.\ 1-kb), the decisive test is to obtain the original
1-kb tumor/normal WGS for neuroblastoma via the dbGaP Data Access Request and
evaluate the predictors \emph{out of sample}. If the prognostic signal returns at
1-kb in NBL, the original claim is genuine but narrow (it requires that
resolution and that biology). If it remains at chance out-of-sample even on the
original data type, the headline $C \approx 0.80$ reflects in-sample evaluation.
Either outcome is a definitive, publishable result---and our open-data panel has
already isolated exactly which variable (resolution $\times$ biology) the answer
turns on.

\subsection{Acknowledged caveats in the open-data pipeline}

In the spirit of honest reproduction, we record the limitations of \emph{our own}
open-data pipeline that bound how strongly the chance-level prognosis can be read
as evidence about the original method. These are properties of the substrate and
construction we used, verified directly from the analysis code, not inferences:

\begin{enumerate}[leftmargin=*]
  \item \textbf{Resolution is intrinsic to the open data, and is the binding
    caveat.} The open TCGA copy-number layer is the ``Masked Copy Number
    Segment'' product (DNAcopy / CBS segmentation of Affymetrix SNP6 arrays). Our
    pipeline bins these segments to 1-Mb genome-wide bins
    (\texttt{BIN = 1{,}000{,}000} in \texttt{f4\_build\_and\_gsvd.py}). This is a
    deliberate, appropriate choice \emph{for segment-level array data}---the SNP6
    segment product does not carry meaningful sub-megabase structure, so finer
    binning would add empty resolution. But it means the $\sim$1000$\times$ gap to
    the paper's 1-kb WGS features is \textbf{inherent to the open data source},
    not a tunable knob. We therefore cannot, on open data, present the method with
    the focal sub-megabase signal it was designed to exploit. This is the single
    most important limitation on interpreting the prognostic null.
  \item \textbf{The matched normal is array-based, not WGS read-depth.} Our
    $Z_2$ (the patient-matched normal genome) is the Blood-Derived-Normal
    aliquot's \emph{own} SNP6 segment profile---a genuine independent genome,
    which is the correct construction in principle. However, it is array-segment
    copy number, whereas the paper's $Z_2$ is WGS read-count depth. The two are
    different measurement modalities of ``the normal genome,'' and any
    modality-specific structure the method relies on would not be present in the
    array substrate.
  \item \textbf{Out-of-sample by cohort, not by the original split.} Our
    concordances are out-of-sample on independent cancers. They are not an
    out-of-sample re-evaluation of the \emph{original} predictors on the
    \emph{original} data type---which only Route~A (the dbGaP 1-kb WGS) can
    provide, and which is the clean comparison to the paper's in-sample
    $C \approx 0.80$.
  \item \textbf{Artifact hygiene.} The per-cohort result files in the F6 panel
    carry a templated \texttt{cohort} field that reads \texttt{"TCGA-BRCA"} in
    every file; the authoritative cancer identity is the filename
    (\texttt{f6\_\{OV,LUAD,SARC,ESCA,STAD,UCEC\}\_results.json}). This is a
    cosmetic logging bug with no effect on the computed numbers, noted so the
    artifacts are read correctly.
\end{enumerate}

\noindent
Taken together, these caveats mean the correct reading of the prognostic null is
narrow and precise: \emph{on segment-resolution open array data, across these
cancers, the unsupervised GSVD predictors do not beat standard baselines out of
sample.} It does \textbf{not} establish that the method fails on its native 1-kb
WGS substrate in copy-number-driven disease---that remains open and is exactly
what Route~A tests. The geometry and two-mechanism reproductions, by contrast,
are robust to all four caveats above (they do not depend on resolution or
modality), which is why those layers reproduce cleanly while prognosis does not.

%====================================================================
\section{Summary}

\begin{center}
\begin{longtable}{p{6.8cm} l p{5cm}}
\toprule
\rowcolor{headgrey}
\textbf{Claim layer} & \textbf{Status} & \textbf{Evidence} \\
\midrule
\endhead
Geometry (tumor-exclusive arraylet) & \textcolor{good}{Reproduces} &
$c_{\text{first}}=1.000$, $c_{\text{last}}\ll 1$ in all 7 native cohorts \\
\addlinespace[2pt]
Two-mechanism independence & \textcolor{good}{Reproduces} &
cross-arraylet $\cos \le 0.05$ in 5/7 cohorts \\
\addlinespace[2pt]
Prognosis, copy-number-driven biology & \textcolor{gap}{Open} &
OV/BRCA tested at 1-Mb; finer resolution untested (Route~A) \\
\addlinespace[2pt]
Prognosis, point-mutation-driven biology & \textcolor{gap}{Did not carry over} &
glioma predictor $C=0.28$--$0.51$ vs.\ age $0.74$ \\
\addlinespace[2pt]
Prognosis, general (open cohorts, 1-Mb) & \textcolor{gap}{Did not carry over} &
all predictor $C \in [0.33,0.56]$, all below Cox baseline \\
\addlinespace[2pt]
Exact original data layer (1-kb WGS) & \textcolor{gap}{Open} &
dbGaP DAR route identified (\texttt{phs000467}, NCI DAC) \\
\bottomrule
\caption{Status of each claim layer.}
\end{longtable}
\end{center}

\paragraph{Bottom line.}
The mathematical core of the method---a GSVD that isolates an orthogonal,
tumor-exclusive copy-number structure from patient-matched tumor and normal
genomes---\textbf{reproduces independently, across ten cohorts and at up to
$11\times$ the original scale.} The strong out-of-sample survival-prediction
result did not transfer to any open cohort we could test; the most parsimonious
explanation is that genome-wide copy number at 1-Mb resolution carries prognosis
only in specific (copy-number-driven) disease biology, and that the original
headline reflects in-sample evaluation on a copy-number-driven cancer at finer
resolution. The decisive next experiment is a clean out-of-sample test on the
original 1-kb data type, obtainable directly via a dbGaP Data Access Request.

This is replication functioning exactly as intended: confirming the robust,
general core of a method, and precisely delimiting the conditions under which its
strongest claim holds---on independent, open data, with pre-registered
expectations and every figure traceable to a committed script.

%====================================================================
\section*{Reproducibility statement}
\addcontentsline{toc}{section}{Reproducibility statement}
All figures derive from committed analysis scripts on open GDC / TCGA data
(\texttt{/data/stevens/alter-pathB/}), using an independent reimplementation of
the Alter--Brown--Botstein / Van Loan GSVD (\texttt{gsvd\_reference.py}).
Per-cohort results: \texttt{f1\_baseline\_table.json}, \texttt{f2\_results.json},
\texttt{f2b\_two\_omic.json}, \texttt{f3\_results.json},
\texttt{f4\_glioma\_results.json}, \texttt{f5\_brca\_results.json}, and
\texttt{f6\_\{OV,LUAD,SARC,ESCA,STAD,UCEC\}\_results.json}. No controlled-access
data was used in any result reported here. Compute: uicgpu (linear-algebra
workload, CPU). Survival engine: \texttt{lifelines}.

\end{document}
